\chapter{Theory of specialization of the~fundamental~group}
\label{X}
% original p. 261
In the present expos\'e, we confine ourselves to the study of the
fundamental group of the geometric fibers in a
\emph{proper} morphism, \ie of the fundamental group of a variable
\emph{proper} algebraic scheme. In a later
expos\'e, we shall generalize the technique employed
to \emph{tamely ramified} \'etale coverings
``at infinity''. This will give us for example a solution of the
``Three-point problem'' in the case of Galois
coverings of order prime to the characteristic (\ie a
determination of the Galois coverings of the line
$\PP_k^1$, ramified at most at three given points and
tamely ramified at these points), and of its obvious
variants.

\section[The homotopy exact sequence for a proper morphism]
{The homotopy exact sequence for a proper and separable morphism}
\label{X.1}

\begin{definition}
\label{X.1.1}
A prescheme $X$ over a field $k$ is said to be
\emph{separable}, or \emph{separable over~$k$}, if for every
extension $K$ of~$k$, $X\otimes_k K$ is reduced. If $f\colon X\to
Y$ is a morphism of preschemes, one says that $f$ is
\emph{separable}, or that $X$ \emph{is separable over~$Y$}, if
$X$ is flat over~$Y$ and if for every $y\in Y$, the fiber $X\otimes_Y
\kres(y)$ is separable over~$\kres(y)$.
\end{definition}

If $X$ is a prescheme over a field~$k$, to say that it is
separable also means that it is \emph{reduced}, and that the
fields~$\kres(x)$, for $x$ a generic point of an irreducible
component of~$X$, are separable extensions of~$k$. If
$k$ is perfect, it therefore amounts to the same to say that $X$ is
reduced. Let us note that if $X$ is separable over~$Y$, then for
every change of base $Y'\to Y$, $X'=X\times_Y Y'$ is separable
over~$Y'$. One can also prove, under suitable finiteness
hypotheses, that the composite of two separable
morphisms is
% original p. 262
a separable morphism. We shall need this only in the
following form: \emph{if $X$ is separable over~$Y$, and $X'$
\'etale over~$X$, $X'$ is separable over~$Y$.} This is indeed
an immediate consequence of the definitions
and~I~\Ref{I.9.2}. Moreover, the hypothesis ``separable
morphism'' will serve us through the following
proposition:
\begin{proposition}
\label{X.1.2}
Let $f\colon X\to Y$ be a proper and separable morphism, with $Y$
locally noetherian, and consider its \emph{Stein}
factorization $X\lto{f'}Y'\dpl\to Y$, (where
$f_*'(\cal{O}_X)=\cal{O}_{Y'}$, $Y'$ being finite over~$Y$ and
isomorphic to the spectrum of the algebra $f_*(\cal{O}_X)$). Then $Y'$
is an \emph{\'etale covering} of~$Y$.
\end{proposition}

This proposition will appear in (EGA~III 7)\footnote{\Cf EGA~III
7.8.10 (i)}. Let us indicate the principle of the proof. One easily
reduces to the case where $Y$ is the spectrum of a complete local
ring~$A$, and, making moreover a suitable finite flat
extension of the latter (corresponding to a suitable
residual extension), one can assume that the connected
components of the fiber of the closed point $y$ are
geometrically connected, which also means that
$\H^0(X_y,\cal{O}_{X_y})$ decomposes into a product of fields
identical to $k=\kres(y)$. Assuming then $X$ connected, which is
permissible, one will have $\H^0(X_y,\cal{O}_{X_y})=k$, hence the homomorphism
$A\to \H^0(X_y,\cal{O}_{X_y})$ is \emph{surjective}. One concludes from this by
a general proposition (of K\"unneth type) that
$f_*(\cal{O}_X)$ is defined by a module $B$ over~$A$ which is
free over~$A$, and that $B/\goth{m} B\to \H^0(X_y,\cal{O}_{X_y})=k$ is
bijective. Hence in the present case $B$ is an \'etale algebra
over~$A$, which completes the proof.

\begin{theorem}
\label{X.1.3}
Let $f\colon X\to Y$ be a proper and separable morphism, with $Y$
locally noetherian and connected, and assume
$f_*(\cal{O}_X)=\cal{O}_Y$ (which implies that the fibers of~$X$ over~$Y$ are geometrically connected, and conversely
thanks to~\Ref{X.1.2}). Let $y$ be a point of~$Y$,
$\overline{\kres(y)}$ an algebraic closure of~$\kres(y)$,
$\overline X_y =X_y\otimes_{\kres(y)} \overline{\kres(y)}$. Finally, let $X'$ be a \emph{connected} \'etale covering of~$X$, and $\overline
X_y' =X_y'\otimes_{\kres(y)} \overline{\kres(y)}$. For there to exist an
\'etale covering $Y'$ of~$Y$ and an $X$-isomorphism
% original p. 263
$X'\isomto X\times_Y Y'$, it is necessary and sufficient that $\overline{X'}_y$
admit a section over~$\overline X_y$.
\end{theorem}

Setting $Y'=\Spec(h_*(\cal{O}_{X'}))$ (where $h\colon X'\to Y$ is the
composite morphism $X'\to X\to Y$), it suffices to prove that the
canonical $Y$-morphism
$$
X'\to X\times_Y Y'
$$
is an \emph{isomorphism}, and that $Y'$ is \'etale over~$Y$. Now we
already know by~\Ref{X.1.2} that $Y'$ is \'etale over~$Y$,
hence $X\times_Y Y'$ is \'etale over~$X$, hence the morphism $X'\to
X\times_Y Y'$ is also \'etale (I~\Ref{I.4.8}). Moreover,
$Y'$ is connected as the image of~$X'$ which is so, hence $X\times_Y Y'$
is connected since $X$ has connected fibers over~$Y$
(IX~\Ref{IX.3.4} and V~\Ref{V.6.9}~(iii)). Hence to prove that
$X'\to X\times_Y Y'$ is an isomorphism, it suffices to see that its
degree of projection at \emph{one} point of~$X\times_Y Y'$ is
equal to~1. Now this follows easily from the hypothesis
that $\overline{X'}_y$ admits a section over~$\overline X_y$, either
by use of IX~\Ref{IX.6.6}, or more simply by noting that it
suffices to prove the existence of such a point in $X\times_Y Y'$
after change of base
$\Spec(\overline{\kres(y)})\to Y$, where this
is evident. This completes the proof of~\Ref{X.1.3}.

Taking into account IX~\Ref{IX.3.4} and the dictionary V~\Ref{V.6.9} and
V~\Ref{V.6.11}, one can put~\Ref{X.1.3} in the following
equivalent form

\begin{corollary}
\label{X.1.4}
With the preceding notation for $f\colon X\to Y$ and
$\overline X_y$, let $\bar{a}$ be a geometric point of
$\overline X_y$, $a$ its image in $X$ and $b$ its image
in~$Y$. Then the following sequence of group homomorphisms is
\emph{exact}:
$$
\pi_1(\overline X_y,\bar{a})\to \pi_1(X,a)\to \pi_1(Y,b)\to e.
$$
\end{corollary}

\begin{remarks}
\label{X.1.5}
One will note that the proof of~\Ref{X.1.3} makes essential
use of~\Ref{X.1.2} and thereby of the ``first
comparison theorem'' in formal algebraic
geometry. On the other hand, the descent theory of
Expos\'e~\Ref{IX}
% original p. 264
has intervened only through IX~\Ref{IX.3.4}, of which
a direct proof in the case of a \emph{proper} morphism
$f\colon X\to Y$ such that $f_*(\cal{O}_X)=\cal{O}_Y$ is easy. Indeed,
let $Y'$ be \'etale over~$Y$ and assume that $X' = X\times_Y Y'$
is the disjoint sum of two nonempty open sets; let us prove that the
same holds for~$Y'$. Indeed, one will have $Y'=\Spec(\cal{A})$, hence
$X'=\Spec(\cal{B})$ with $\cal{B}=\cal{A}\otimes_{\cal{O}_Y}
\cal{O}_X$, and the decomposition of~$X'$ as a direct sum
corresponds to a decomposition of~$\cal{B}$ as a product of two
nonzero Algebras $\cal{B}_1$ and $\cal{B}_2$. Since
$f_*(\cal{O}_X)=\cal{O}_Y$ one easily concludes
$f_*(\cal{B})=\cal{A}$, hence $\cal{A}$ will be the sum of two Algebras
(likewise nonzero, since their unit sections are
nonzero) $f_*(\cal{B}_1)$ and $f_*(\cal{B}_2)$, qed.
\end{remarks}

\subsection{}
\label{X.I.6}
Let us again assume that $f$ is proper and separable, but let us no longer make
any hypothesis on~$f_*(\cal{O}_X)$, which will correspond to a
well-determined \'etale covering $Y'$ of~$Y$,
moreover pointed over~$b$ by the image $b'$
of~$a$. Applying then \Ref{X.1.4} to the canonical morphism $X\to Y'$,
and assuming $f$ surjective, the exact sequence~\Ref{X.1.4} is
replaced by the following one, analogue of the homotopy exact sequence
of fiber spaces in algebraic topology:
$$
\pi_1(\overline X_y,\overline a)\to \pi_1(X,a)\to \pi_1(Y,b)\to
\pi_0(\overline X_y,\overline a)\to \pi_0(X,a)\to \pi_0(Y,b)\to e.
$$
Of course, in~\Ref{X.1.4} one cannot in general
assert that the homomorphism
$\pi_1(\overline X_y,\overline a)\to \pi_1(X,a)$
is injective; in algebraic topology, its kernel is the image of
$\pi_2(Y,b)$, and there would be occasion in algebraic geometry
as well to introduce homotopy groups in all
dimensions, and the complete homotopy exact sequence for a
proper morphism satisfying suitable hypotheses (for
example that of being a smooth
morphism).
One does not have at present
any result in this direction, with the exception of a
reasonable (if not definitive) definition of the higher homotopy
groups.

\begin{corollary}
\label{X.1.7}
Let $k$ be an algebraically closed field, $X$ and $Y$ two
connected preschemes over~$k$; one assumes $X$ proper over~$k$ and
$Y$ locally
% original p. 265
noetherian. Let $a$ be a geometric point of~$X$, $b$
a geometric point of~$Y$ with values in the same
algebraically closed extension $K$ of~$k$; consider the
geometric point $c=(a,b)$ of~$X\times_k Y$, and
the homomorphism $\pi_1(X\times_k Y,c)\to \pi_1(X,a)\times\pi_1(Y,b)$
deduced from the homomorphisms on the fundamental groups
associated with the two projections $X\times_k Y\to X$ and $X\times_k
Y\to Y$. The preceding homomorphism is an \emph{isomorphism}.
\end{corollary}

Let us first assume $K=k$. Let us set $Z=X\times_k Y$, consider the
projection $f\colon Z\to Y$ and the locality $y$ of the
geometric point $b$ of~$Y$, and let us apply to the situation the
result~\Ref{X.1.4}. One will note for this that, at the cost of passing
to $X_\red$ (which does not change the fundamental groups
under consideration), one can already assume $X$ reduced, hence
separable over~$k$, hence~$Z$ is separable over~$k$, and
evidently with geometrically connected fibers (since~$X$ is connected). The geometric fiber of~$Z$ at $b$ is
canonically isomorphic to $X\otimes_k K=X$. On the other hand, since the
composite of the morphisms $X\to Z\to X$ is the identity, one finds
that $\pi_1(X,a)\to \pi_1(Z,c)$ is injective, and 1.4 gives us an
exact sequence:
$$
e\to \pi_1(X,a)\to \pi_1(Z,c)\to \pi_1(Y,b)\to e.
$$
On the other hand, one has the canonical exact sequence:
$$
e\to \pi_1(X,a)\to \pi_1(X,a)\times \pi_1(Y,b)\to \pi_1(Y,b)\to e,
$$
where the two homomorphisms written are the canonical injection
and the canonical projection. Finally, one has a homomorphism from the
first exact sequence into the second, by means of the
identity morphisms on the extreme terms, and the canonical
homomorphism $\pi_1(Z,c)\to \pi_1(X,a)\times\pi_1(Y,b)$ for the middle
terms. The commutativity of the diagram thus obtained is
verified trivially. Since the homomorphisms on the extreme
terms are isomorphisms, the same holds for the middle
terms, which proves \Ref{X.1.7} in this case.

When one no longer assumes $K=k$, one finds only an isomorphism
$$
\pi_1(Z,c)\to \pi_1(X\otimes_k K,a)\times \pi_1(Y,b),
$$
% original p. 266
and~\Ref{X.1.7} is then equivalent to the following particular case:

\begin{corollary}
\label{X.1.8}
Let $X$ be a proper and connected scheme over an
algebraically closed field~$k$, $k'$ an algebraically
closed extension of~$k$, $a'$ a geometric point of~$X\otimes_k k'$ and
$a$ its image in~$X$. Then the canonical homomorphism
$\pi_1(X\otimes_k k',a')\to \pi_1(X,a)$ is an \emph{isomorphism.}
\end{corollary}

The fact that this homomorphism is surjective is equivalent to saying
that if $X'$ is a connected \'etale covering of~$X$, then
$X'\otimes_k k'$ is also connected, and follows
immediately from the fact that $k$ is algebraically closed; it is also
a particular case
of~IX~\Ref{IX.3.4}.
The hypothesis of properness on~$X$ has not yet been used. That
said, to say that the homomorphism under consideration is injective also means
this: \emph{every \'etale covering of~$X\otimes_k k'$ is
isomorphic to the inverse image of an \'etale covering of~$X$.}
It is essentially a matter of sorites that one can find a
sub-$k$-algebra $A$ of
$k'$,
of finite type over~$k$, and an
\'etale covering of~$X\otimes_k A$ whose inverse image on~$X\otimes_k k'$ is isomorphic to the given covering. Let then
$Y=\Spec(A)$, which is an integral $k$-scheme of finite type, hence
has rational points over~$k$. Let us then apply \Ref{X.1.7} to the
fundamental group of~$X\times Y$ at a point $(a,b)$ rational
over~$k$: one finds that every connected \'etale covering of
$X\times Y$ is isomorphic to a quotient of a covering
$X'\times Y'$, where $X'$, $Y'$ are Galois \'etale
coverings of~$X$ and $Y$ with groups~$G$, $G'$, by a subgroup $H$
of~$G\times G'$. This implies that the inverse image of this
covering of~$X\times Y$ on~$X\times Y'$ is isomorphic to a
covering of the form $X_1'\times Y'$, where $X_1'$ is an
\'etale covering of~$X$. If then $L$ is the field of
functions of~$Y$, equal to the field of fractions of~$A$ in~$k'$, the
\'etale covering of~$X\otimes_k L$ induced by the
given covering of~$X\times_k Y$ is such that there exists a
finite separable extension $L'$ of~$L$ such that the inverse image
of the said covering on~$X\otimes_k L'$ is isomorphic to
$X_1'\otimes_k L'$. Now, $k'$ being algebraically closed, one can
assume that the extension $L'$ of~$L$ is contained in~$k'$. This
proves that the given \'etale covering of~$X\otimes_k k'$
is isomorphic to $X_1'\otimes_k k'$, qed.

The explicit form noted in passing for the \'etale
coverings of a product
% original p. 267
$X\times_k Y$ immediately implies the following result:

\begin{corollary}
\label{X.1.9}
Let $k$ be an algebraically closed field, $X$ and $Y$ two
locally noetherian preschemes over~$k$, $Z=X\times_k
Y$ their product, $Z'$ an \'etale covering of~$Z$. For every
point $y\in Y$ rational over~$k$, let $i_y \colon \Spec(k)\to Y$ be the
canonically associated morphism, $j_y = \id_X\times_k i_y$ the
corresponding morphism $X\to Z$. Finally, let $X_y'$ be the \'etale
covering of~$X$ inverse image of~$Z'$ by~$j_y$. One assumes $Y$
connected, and $X$ or~$Y$ proper over~$k$. Then the coverings
$X_y'$ of~$X$ are all isomorphic.
\end{corollary}

Figuratively, one can say that \emph{a family of
\'etale coverings of~$X$, parametrized by a
connected prescheme~$Y$, is constant if $X$ or the
parameter prescheme~$Y$ is proper over~$k$.}

\begin{remarks}
\label{X.1.10}
Corollaries~\Ref{X.1.7} to~\Ref{X.1.9} are due to
Lang-Serre~\cite{X.2} in the case of normal algebraic schemes.
(Their work was the initial motivation for the theory
of the fundamental group developed in this Seminar). As
these authors have remarked, these results become false
when one drops the hypothesis of properness, at least in
characteristic $p>0$. Taking for example for $X$ the affine
line $X=\Spec(k[t])$, it is not difficult to see that the
coverings of~$X$, parametrized by the affine line
$Y=\Spec(k[s])$, defined by the equations
$$
x^p - x = st,
$$
are \'etale and pairwise non-isomorphic. This invalidates~\Ref{X.1.9}
and a fortiori~\Ref{X.1.7}, and one sees
similarly that if $s$ is considered as an element
transcendental over~$k$ in an algebraically closed extension $K$
of~$k$, one finds an \'etale covering $X'$ of~$X\otimes_kK$
which does not come from an \'etale covering of~$X$.
\end{remarks}
